% Generated draft: requires mathematical and source review.
\documentclass[11pt]{article}
\usepackage[margin=25mm]{geometry}
\usepackage{fontspec}
\setmainfont{Latin Modern Roman}
\usepackage{amsmath,amssymb,mathtools,unicode-math}
\setmathfont{Latin Modern Math}
\usepackage{enumitem,needspace,xurl,hyperref}
\setlength{\parindent}{0pt}
\setlength{\parskip}{5pt}
\setlength{\emergencystretch}{3em}
\newcommand{\N}{\mathbb{N}}
\newcommand{\Z}{\mathbb{Z}}
\newcommand{\Q}{\mathbb{Q}}
\newcommand{\R}{\mathbb{R}}
\newcommand{\C}{\mathbb{C}}
\newcommand{\F}{\mathbb{F}}
\newcommand{\A}{\mathbb{A}}
\newcommand{\PP}{\mathbb{P}}
\newcommand{\E}{\mathbb{E}}
\newcommand{\eps}{\varepsilon}
\newcommand{\dd}{\,\mathrm{d}}
\begin{document}
\section*{UnsolvedMath ID: 30002139}
\textbf{Source problem code:} catalogue entry
\subsection{Definitions and mathematical statement}
The supplied record asks about the geometry of horospheres for lifted metrics, but omits hypotheses and quantifiers needed for a unique precise formulation. The following definitions describe a conventional setting; completeness, the relationship between the underlying manifolds, and the choice of metrics on horospheres remain qualifications requiring review.

Let \((X,h)\) be a complete simply connected Riemannian manifold of nonpositive sectional curvature. For a unit-speed geodesic ray \(\gamma:[0,\infty)\to X\), define its Busemann function and associated horospheres by
\[
 b_\gamma(x)=\lim_{t\to\infty}\bigl(d_h(x,\gamma(t))-t\bigr),
 \qquad H_{\gamma,s}=b_\gamma^{-1}(s),\quad s\in\R.
\]
When a horosphere is equipped with its intrinsic Riemannian distance, this distance is
\[
 d_H(x,y)=\inf\{\operatorname{Length}_h(c):c\text{ is a piecewise smooth curve in }H\text{ joining }x\text{ to }y\}.
\]
The restricted ambient distance \(d_h|_{H\times H}\) is a different possible convention. The supplied record does not specify which is intended.

A map \(F:(Y,d_Y)\to(Z,d_Z)\) is a quasi-isometry if there exist constants \(L\geq1\) and \(A\geq0\) such that, for every \(y,y'\in Y\),
\[
 L^{-1}d_Y(y,y')-A\leq d_Z(F(y),F(y'))\leq Ld_Y(y,y')+A,
\]
and every point of \(Z\) has distance at most \(A\) from \(F(Y)\).

\textbf{Dimension-three question, with unresolved scope.} Let \(\widetilde g_0\) denote the lifted hyperbolic metric mentioned in the source, and let \(\widetilde g\) denote the lift of another negatively curved metric. In dimension \(3\), the source asks whether their horospheres are quasi-isometric:
\[
 (H_0,d_{H_0})\simeq_{\mathrm{qi}}(H,d_H).
\]
Here \(H_0\) and \(H\) are horospheres for the respective lifted metrics. The source does not specify whether these metrics are on the same manifold, whether that manifold is closed or complete, whether curvature is pinched, or whether the comparison concerns every pair of horospheres, some pair, or pairs corresponding under a specified boundary identification. Thus the displayed question retains unresolved hypotheses and quantifiers; it is not a fully specified universal assertion.

\textbf{General question.} Given a quasi-isometry between universal covers carrying negatively curved metrics, determine what additional hypotheses, if any, yield a comparison of their horospheres. The source does not specify whether the desired conclusion is the existence of a quasi-isometry between horospheres or a conclusion concerning the given ambient quasi-isometry and its images of horospheres.

\textbf{Volume-growth question.} The source additionally asks under what conditions its lower and upper polynomial volume-growth exponents \(\alpha\) and \(\beta\) coincide. It supplies no definitions of these exponents. One possible convention, if intrinsic Riemannian balls and volume are intended, would be to put
\[
 V_{H,p}(r)=\operatorname{vol}_H\bigl(B_{d_H}(p,r)\bigr),\qquad
 \underline\delta(H,p)=\liminf_{r\to\infty}\frac{\log V_{H,p}(r)}{\log r},\qquad
 \overline\delta(H,p)=\limsup_{r\to\infty}\frac{\log V_{H,p}(r)}{\log r}.
\]
These formulas require the indicated volumes to be finite and positive for sufficiently large \(r\). Identifying \(\alpha=\underline\delta(H,p)\) and \(\beta=\overline\delta(H,p)\) would be an additional interpretation, not a definition established by the supplied record. Whether these are the intended quantities, which horospheres and base points are quantified over, and whether finite polynomial growth is assumed must be resolved before the equality question can be stated precisely.
\subsection{Short English statement}
The source asks whether horospheres for a lifted hyperbolic metric and another lifted negatively curved metric are quasi-isometric in dimension \(3\), what analogous conclusions follow from quasi-isometries of universal covers in higher dimensions, and when lower and upper polynomial volume-growth exponents agree. The manifold hypotheses, horosphere metrics, comparison quantifiers, and exponent definitions are missing, so this draft records the questions with explicit unresolved qualifications.
\subsection{Sources}
\begin{itemize}
\item Supplied local UnsolvedMath registry record, permanent ID \(30002139\), title: Horosphere Geometry Under Quasi-Isometries of Universal Covers; exact source problem code: \texttt{catalogue entry}. Collection: Oberwolfach Reports Open Problems. The proposed year is recorded as \(2012\); no proposer is named.
\item Supplied report citation: \emph{Geometric Group Theory, Hyperbolic Dynamics and Symplectic Geometry} (2013), Oberwolfach Reports, pp.\ 2139--2203, DOI \texttt{10.4171/owr/2012/35}. Also listed among the supplied literature sources as Oberwolfach Report 35: \url{https://doi.org/10.4171/owr/2012/35}. The supplied background describes these reports as records of mathematical workshops and their open-problem discussions.
\item Supplied literature reference: Besson--Courtois--Hersonsky, horospherical rigidity, \url{https://arxiv.org/abs/2411.15093}.
\item Supplied literature reference: Rigidity of flat holonomies, \url{https://doi.org/10.1017/etds.2024.58}.
\item Literature assessment dated exactly \texttt{2026-08-21}: the supplied curation labels the problem open and states that the dimension-three hyperbolic comparison and the proposed higher-dimensional generalization remain unresolved in the stated generality. This is supplied dated curation, not an independently verified assessment. The background explicitly requests independent verification of time-sensitive status and cited claims.
\item The registry labels the statement \texttt{corrected\_verified} and reports that its mathematical content was checked against the official Oberwolfach report, with notation, context, and question form normalized without changing the source-level claim. This verification claim is supplied provenance; the report and literature references have not been independently consulted for this draft.
\end{itemize}
\end{document}
